\subsection{关于绝对值的预备知识}

设 K 是一个域（交换或非交换）。回顾（《一般拓扑学》第九章 §3 第 2 节定义 2），K 上的绝对值是指 K 到 \(R_{+}\) 的满足下列公理的任一映射 f：

\((\mathrm{VA}_{\mathbf{I}})\) 关系 \(f(x) = 0\) 等价于 \(x = 0\) 。

\((\mathrm{VA}_{\mathrm{II}}) f(xy) = f(x)f(y)\) 对一切 \(x, y\) （在 \(\mathbf{K}\) 中）成立。

\((\mathrm{VA}_{\mathrm{III}}) f(x + y) \leqslant f(x) + f(y)\) 对一切 \(x, y\) （在 \(\mathbf{K}\) 中）成立。

由 \((\mathrm{VA}_{\mathbf{I}})\) 与 \((\mathrm{VA}_{\mathbf{II}})\) 可得 \(f(1) = 1, f(-1) = 1\) 以及

\[
f (x ^ {- 1}) = \frac {1}{f (x)}
\]

对 \(x \neq 0\) 成立。

对映射 \(f\) （由 \(\mathbf{K}\) 到 \(\mathbf{R}_{+}\) ）与实数 \(\mathbf{A} > 0\) ，用 \((\mathrm{U}_{\mathbf{A}})\) 表示关系

\[
f (x + y) \leqslant \mathrm{A}. \sup (f (x), f (y)) \quad \text {   for   all   } x, y \text {   in   } \mathrm{K}.
\]

我们用 \(\mathcal{V}(\mathbf{K})\) 表示由映射 \(f\) （从 \(\mathbf{K}\) 到 \(\mathbf{R}_{+}\) ）组成的集合，这些映射满足 \((\mathrm{VA}_{\mathbf{I}})\) 与 \((\mathrm{VA}_{\mathbf{II}})\) ，且存在 \(\mathbf{A} > 0\) （依赖于 \(f\) ）使 \((\mathbf{U}_{\mathbf{A}})\) 成立。

注意，若 \(f \in \mathcal{V}(\mathbf{K})\) ，则把 \(x = 1, y = 0\) 代入 \((\mathrm{U}_{\mathbf{A}})\) ，

\[
1 = f (1) \leqslant A. \sup (f (1), f (0)) = A.
\]

命题 1. 映射 \(f\) （从 \(\mathbf{K}\) 到 \(\mathbf{R}_+\) ）满足 \((\mathrm{VA}_{\mathbf{I}})\) 与 \((\mathrm{VA}_{\mathbf{II}})\) 而属于 \(\mathcal{V}(\mathbf{K})\) ，当且仅当 \(f(1 + x)\) 在由这样的 \(x \in \mathbf{K}\) （满足 \(f(x) \leqslant 1\) ）组成的集合上有界。

若 \(f\) 满足 \((\mathrm{U}_{\mathbf{A}})\) ，则 \(f(1 + x) \leqslant \mathbf{A}\) 当 \(f(x) \leqslant 1\) 时成立。反之，设 \(f(x + 1) \leqslant \mathbf{A}\) 对这样的 \(x \in \mathbf{K}\) （满足 \(f(x) \leqslant 1\) ）成立（这蕴含 \(\mathbf{A} \geqslant f(1) = 1\) ）；那么，若 \(x = 0\) 或 \(y = 0\) ，条件 \((\mathrm{U}_{\mathbf{A}})\) 成立；另一方面，若 \(x \neq 0\) 且 \(y \neq 0\) ，不妨设 \(f(y) \leqslant f(x)\) ，于是由 \((\mathrm{VA}_{\mathrm{II}}), f(yx^{-1}) \leqslant 1\) ，因此 \(f(1 + yx^{-1}) \leqslant \mathbf{A}\) ，由 \((\mathrm{VA}_{\mathrm{II}}), f(x + y)f(x)^{-1} \leqslant \mathbf{A}\) ；从而

\[
f (x + y) \leqslant A f (x) \leqslant A. \sup (f (x), f (y)).
\]

若 \(f\) 是 \(\mathbf{K}\) 上的绝对值，则 \(f(n,1) \leqslant n\) 对整数 \(n > 0\) 从 \((\mathrm{VA}_{\mathrm{III}})\) 出发用归纳法可得；反之：

命题 2. 设 \(f\) 是 \(\mathbf{K}\) 到 \(\mathbf{R}_+\) 的一个映射，且属于 \(\mathcal{V}(\mathbf{K})\) ；若存在 \(\mathbf{C} > 0\) 使得 \(f(n,1) \leqslant \mathbf{C}.n\) 对每个整数 \(n > 0\) 成立，则 \(f\) 是 \(\mathbf{K}\) 上的绝对值。

对 r > 0 作归纳，由 \((\mathbf{U}_{\mathbf{A}})\) 推出关系

\[
f (x _ {1} + x _ {2} + \dots + x _ {2 ^ {r}}) \leqslant A ^ {r} \sup _ {1 \leqslant i \leqslant 2 ^ {r}} f (x _ {i})\tag{1}
\]

对每个族 \((x_{i})\) ——由 \(2^{r}\) 个 \(\mathbf{K}\) 的元素组成——成立。令 \(n = 2^{r} - 1\) ；对一切 \(x \in \mathbf{K}\) ，由 (1) 推出

\[
(f (1 + x)) ^ {n} = f ((1 + x) ^ {n}) = f \left(\sum_ {i = 0} ^ {n} \binom {n} {i} x ^ {i}\right) \leqslant A ^ {r} \sup \left(f \left(\binom {n} {i}\right) (f (x)) ^ {i}\right)
\]

\[
\leqslant \mathrm{CA} ^ {r} \sum_ {i = 0} ^ {n} \binom {n} {i} (f (x)) ^ {i} = \mathrm{CA} ^ {r} (1 + f (x)) ^ {n}
\]

因为 \(f\left(\binom{n}{i}\right) \leqslant C\binom{n}{i}\) ；因此

\[
f (1 + x) \leqslant \mathrm{C} ^ {1 / n} \mathrm{A} ^ {r / n} (1 + f (x)).
\]

令 r 趋于 \(+\infty\) ，得到 \(f(1+x)\leqslant1+f(x)\) 对一切 \(x\in K\) 成立；把 x 换成 \(xy^{-1}\) （其中 \(y\neq0\) ）应用这个不等式，并注意到 \((\mathrm{VA}_{\mathrm{II}})\) ，就得到关系 \((\mathrm{VA}_{\mathrm{III}})\) ，这证明了命题。

推论 1. 由 K 到 \(R_{+}\) 的映射 f 是绝对值，当且仅当它满足条件 (VA \(_{I}\) )、(VA \(_{II}\) ) 与 (U \(_{2}\) )。

这是必要的，因为 \((\mathrm{VA}_{\mathrm{II}})\) 蕴含

\[
f (x + y) \leqslant f (x) + f (y) \leqslant 2 \sup (f (x), f (y)).
\]

反之，设 \(f\) 满足 \((\mathrm{VA}_{\mathrm{I}}), (\mathrm{VA}_{\mathrm{II}})\) 与 \((\mathrm{U}_2)\) ；对每个整数 \(n > 0\) ，设 \(r\) 是使 \(2^r \geqslant n\) 的最小整数；若在 (1) 中把 A 换成 2，把 \(x_i\) 中指标 \(i \leqslant n\) 的换成 1，把 \(x_i\) 中指标 \(i > n\) 的换成 0，就得到

\[
f (n. 1) \leqslant 2 ^ {r} <   2 n;
\]

于是可以取 \(\mathbf{C} = 2\) 应用命题 2，因而 \(f\) 是绝对值。

推论 2. 映射 \(f\) （由 \(\mathbf{K}\) 到 \(\mathbf{R}_+\) ）属于 \(\mathcal{V}(\mathbf{K})\) ，当且仅当它形如 \(g^t\) ，其中 \(t > 0\) 且 \(g\) 是 \(\mathbf{K}\) 上的绝对值。

说 \(f\) 满足 \((\mathbf{U}_{\mathbf{A}})\) 等价于说 \(f^s\) 满足 \((\mathbf{U}_{\mathbf{A}^s})\) ；由于存在 \(s > 0\) 使 \(\mathbf{A}^s \leqslant 2\) ，推论 1 表明对这样的值，\(s, f^s\) 是绝对值。

